\chapter{Хавсралт}

\section{Системийн тохиргооны жишээ}

Туршилтын орчинд дараах үндсэн тохиргоог хэрэглэсэн: Python орчин, Firefox WebDriver, өгөгдлийн сангийн холболт, API токен, скрапингийн delay параметрүүд. Эдгээр тохиргоо нь туршилтыг дахин гүйцэтгэхэд шаардлагатай суурь нөхцөлийг бүрдүүлнэ.

\section{API Endpoint жишээ}

Системд ашигласан гол endpoint-уудыг дараах байдлаар бүлэглэж болно.

\begin{table}[H]
\centering
\caption{Хавсралтын endpoint жагсаалт}
\small
\renewcommand{\arraystretch}{1.2}
\begin{adjustbox}{max width=\textwidth}
\begin{tabular}{|p{4.2cm}|p{9.2cm}|}
\hline
\rowcolor{gray!20}\textbf{Endpoint} & \textbf{Зориулалт} \\ \hline
\texttt{/health} & Backend болон dependency status шалгах \\ \hline
\texttt{/api/v1/posts}, \texttt{/api/v1/stats} & Dashboard-ийн үндсэн өгөгдөл унших \\ \hline
\texttt{/api/v1/pages} & Scraper target list унших, хадгалах \\ \hline
\texttt{/api/v1/scraper/run},\newline\texttt{/api/v1/scraper/}\newline\texttt{status},\newline\texttt{/api/v1/scraper/logs} & Scraper ажиллуулах, төлөв болон log харах \\ \hline
\texttt{/api/v1/sentiment}, \texttt{/api/v1/topics}, \texttt{/api/v1/trends} & Analysis workflow ажиллуулах \\ \hline
\texttt{/gemini/analyze}, \texttt{/gemini/ask} & AI тайлан болон асуулт-хариулт үүсгэх \\ \hline
\texttt{/export/sheets} & Google Sheets export хийх \\ \hline
\end{tabular}
\end{adjustbox}
\normalsize
\end{table}

\textbf{API Response JSON жишээ (FastAPI)}:
\begin{verbatim}
{
  "status": "success",
  "data": {
    "total_posts": 150,
    "total_comments": 432,
    "average_sentiment": 0.45,
    "recent_posts": [
      {
        "post_id": "123456789",
        "page_name": "TechNews Mongolia",
        "content": "Шинэ AI технологийн талаарх дэлгэрэнгүй мэдээлэл...",
        "likes": 120,
        "comments": 45
      }
    ]
  }
}
\end{verbatim}

\section{Өгөгдлийн сангийн схем}

Өгөгдлийн сангийн үндсэн схем нь \texttt{FacebookPost}, \texttt{PostComment}, \texttt{PostImage}, \texttt{AnalysisResult} хүснэгтүүд дээр төвлөрсөн. \texttt{FacebookPost} нь үндсэн контент болон engagement metric, \texttt{PostComment} нь сэтгэгдлийн өгөгдөл, \texttt{PostImage} нь дагалдах медиа мэдээлэл, \texttt{AnalysisResult} нь sentiment/topic/engagement шинжилгээний JSON үр дүнг хадгална.
